\kbvchapter{ಪಾಠ ೨೧}{ಛಂದಸ್ಸು}

\kbvsechead{ಪದ್ಯಗಳನ್ನು ರಚಿಸುವುದು ಎಲ್ಲರಿಗೂ ಸಾಧ್ಯವಾಗಲಾರದು :}

\indent ಇದೊಂದು ಹುಟ್ಟು ಕಲೆ, ಈ ಕಲೆ ಕರಗತವಾಗಬೇಕಾದರೆ ಕೆಲವೊಂದು ನಿಯಮಗಳನ್ನು ತಿಳಿದು ಪಾಲಿಸಬೇಕು, ಹೀಗೆ ಪದ್ಯರಚನಾಕೌಶಲ್ಯವನ್ನು ತಿಳಿಸುವ ಶಾಸ್ತ್ರವೇ ಛಂದಶಾಸ್ತ್ರ, ಛಂದಶಾಸ್ತ್ರವನ್ನು ಸಾಕಷ್ಟು ಮಟ್ಟಿಗೆ ಅರ್ಥಮಾಡಿಕೊಳ್ಳದಿದ್ದರೆ ಪದ್ಯಗಳಲ್ಲಿ ಅಡಗಿರುವ ಸೊಬಗನ್ನು ಕಾಣ ಲಾಗುವುದಿಲ್ಲ.\par

\indent “ಛಂದೋಂಬುಧಿ” ಎಂಬ ಹಳಗನ್ನಡ ಗ್ರಂಥದಲ್ಲಿ ಪದ್ಯಗಳ ರಚನೆಗೆ ಸಂಬಂಧಪಟ್ಟ ನಿಯಮಗಳನ್ನು ಮತ್ತು ಪದ್ಯಗಳ ಲಕ್ಷಣಗಳನ್ನು ವಿಸ್ತಾರವಾಗಿ ಹೇಳಲಾಗಿದೆ. ಇದು ಬಹುಮುಖ್ಯವಾದ ಆಧಾರ ಗ್ರಂಥ; ಈ ಗ್ರಂಥದ ಕರ್ತೃ ಕ್ರಿ.ಶ. ಹತ್ತನೆಯ ಶತಮಾನದಲ್ಲಿದ್ದ “ಕವಿರಾಜಹಂಸ” ಬಿರುದಾಂಕಿತ “ನಾಗವರ್ಮ”.\par

\indent ಪದ್ಯಗಳಲ್ಲಿ ನಿರ್ದಿಷ್ಟ ಸಂಖ್ಯೆಯ ಸಾಲುಗಳಿರುವುವು. ಈ ಸಾಲುಗಳಿಗೆ "ಪಾದ”ಗಳು ಅಥವಾ “ಚರಣ”ಗಳು ಎಂದು ಹೆಸರು, ಕನ್ನಡ ಸಾಹಿತ್ಯ ಕ್ಷೇತ್ರದಲ್ಲಿ ಎರಡು, ಮೂರು, ನಾಲ್ಕು ಮತ್ತು ಆರು ಪಾದಗಳ ಪದ್ಯಗಳನ್ನು ಬಹುಮಟ್ಟಿಗೆ ಕಾಣಬಹುದು; ಇವುಗಳನ್ನು ಕ್ರಮವಾಗಿ ದ್ವಿಪದಿ, ತ್ರಿಪದಿ, ಚೌಪದಿ ಮತ್ತು ಷಟ್ಪದಿ ಎನ್ನುವರು.\par

\indent ಯಾವುದೇ ಪದ್ಯದ ಪ್ರತಿಯೊಂದು ಪಾದವೂ ಯತಿ, ಗಣ, ಪ್ರಾಸ ನಿಯಮಗಳಿಗೆ ಅನುಸಾರವಾಗಿರುವುದು.\par

\vspace{4pt plus 4pt minus 1pt}\noindent{\kbvsecfont\fontsize{11pt}{14.5pt}\selectfont\color{kbvaccentdark} (೧)\hspace{0.45em}ಯತಿ:}\hspace{0.5em}ಪದ್ಯವನ್ನು ಓದುವಾಗ ಅವಶ್ಯಕವಾಗಿ ನಿಲ್ಲಿಸುವ ಹಂತವೇ “ಯತಿ” ; ಗೊತ್ತಾದ ಹಂತದಲ್ಲಿ ಓದು ನಿಲ್ಲಿಸಿದಾಗ ಅರ್ಥ ಕೆಡುವಂತಿರಬಾರದು. ಆದರೆ ಕನ್ನಡದಲ್ಲಿ ಕಾವ್ಯರಚನೆ ಮಾಡುವಾಗ ಯತಿನಿಯಮವನ್ನು ಕಟ್ಟುನಿಟ್ಟಾಗಿ ಪಾಲಿಸುವ ಸಂಪ್ರದಾಯ ನಮ್ಮಲ್ಲಿ ಬೆಳೆದುಬಂದಿಲ್ಲ.\par
